\section{代码说明}

研究代码按数据预处理、统计分析、绘图和质量核对组织。所有脚本从项目根目录执行，输入归档由 \nolinkurl{code/00_preprocess.py} 读取。该脚本核对三个 NetCDF 的变量、单位、时间和掩膜，按每月实际天数聚合 DJF 冬季，并把冬季场、备用场核对结果及元数据审计记录分别写入 \nolinkurl{output/processed/}、\nolinkurl{output/csv/} 和 \nolinkurl{logs/}。

\nolinkurl{code/03_global_metrics.py} 与 \nolinkurl{code/04_grid_trends_mk_sen.py} 计算全域年度指标、格点趋势及显著性。\nolinkurl{code/05_region_latitude_analysis.py} 构建代表窗口和纬度带，比较覆盖率、SWE 和气温趋势。\nolinkurl{code/06_temperature_response.py} 计算控制年份后的温度关联、跨冬季滞后与组间差异。\nolinkurl{code/07_threshold_snowline.py} 提取 SCF 阈值等值线、逐经度边界及阈值敏感性。\nolinkurl{code/07b_assimilation_sensitivity.py} 与 \nolinkurl{code/08_swe_tail.py} 检查结果对同化断点和极高 SWE 格点的敏感度。

\nolinkurl{code/08_domain_figure.py}、\nolinkurl{code/08_trend_figures.py}、\nolinkurl{code/08_spatial_figures.py}、\nolinkurl{code/08_snowline_figures.py} 与 \nolinkurl{code/08_synthesis_figures.py} 生成文中图件。\nolinkurl{code/figure_style.py} 统一字体、颜色、地图坐标和导出设置。\nolinkurl{code/09_export_latex_tables.py} 从 CSV 生成正文表格。\nolinkurl{code/10_quality_control.py} 核对年度指标、图件及编译后的报告。\nolinkurl{output/agent_mechanism/assimilation_checks.py} 提供独立的同化断点复核。

运行顺序、依赖版本、随机种子和编译命令见项目根目录的 \nolinkurl{README.md} 与 \nolinkurl{requirements.txt}。脚本读取原始归档并重新生成结果。本文用 AI 辅助编写和核对分析代码、图件及报告文字，统计量由随附脚本计算。
